%% USPSC-Abstract.tex
%\autor{Silva, M. J.}
\begin{resumo}[Abstract]
 \begin{otherlanguage*}{english}
	\begin{flushleft} 
		\setlength{\absparsep}{0pt} % ajusta o espaçamento dos parágrafos do resumo		
 		\SingleSpacing  		\imprimirautorabr~~\textbf{\imprimirtitleabstract}.	\imprimirdata.  \pageref{LastPage} p. 
		%Substitua p. por f. quando utilizar oneside em \documentclass
		%\pageref{LastPage} f.
		\imprimirtipotrabalhoabs~-~\imprimirinstituicao, \imprimirlocal, 	\imprimirdata. 
 	\end{flushleft}
	\OnehalfSpacing 

Understanding how quantum systems equilibrate is a fundamental question with implications for quantum technologies, such as state preparation for quantum computing. However, linear response theory fails to explain the thermal relaxation in the quantum regime and of far-from-equilibrium systems. In this context, recent studies have reported a relaxation asymmetry, where heating is faster than cooling for thermodynamically equidistant initial states, but understanding how quantum resources affect this asymmetry is still an open question. In this work, we investigate the heating and cooling asymmetry in discrete and continuous-variable quantum systems. For the discrete case, we study two interacting qubits coupled to a global bosonic bath, focusing on the effects of initial coherences, steady state correlations generated by the dynamics and local non-Markovianity. For the continuous case, we consider Gaussian states and compare the thermal relaxation asymmetry of displaced and squeezed states. Our results for the discrete case indicate that non-Markovianity may accelerate thermal relaxation, while the initial coherences do not change the equilibration time. In Gaussian systems, applying the displacement operator in the initial states does not affect the asymmetry, but applying the squeezing operator reduces the asymmetry. Additionally, we show an invertion of the asymmetry, making cooling faster than heating, by squeezing the hotter state and displacing the cooler state.

   \vspace{\onelineskip}
 
   \noindent 
   \textbf{Keywords}: asymmetry; relaxation; two-interacting qubits; gaussian systems.
 \end{otherlanguage*}
\end{resumo}
